%% ======================================================================
%% Reviewer Response -- Additional MLLM Benchmark Comparison
%% Addresses reviewer concern: "Limited Benchmark Comparisons"
%% ======================================================================

%% ----------------------------------------------------------
%% TABLE: Extended MLLM benchmark
%% Compares Task 1 (direct classification) and the MBERT-based
%% description-mediated pipeline (Task 2a / 2b) for three new
%% multimodal models (Gemini, Phi-4, Gemma-4-E4B) against the
%% best-performing existing model (GPT-4o mini) reported in the paper.
%% ----------------------------------------------------------

\begin{table*}[!htb]
\begin{adjustwidth}{-2.25in}{0in}
\footnotesize
\caption{%
  Weighted $F$-score (\%, $\pm$95\% CI, 5-fold CV) for Task~1 (direct classification),
  Task~$2_a$ (pre-trained MBERT), and Task~$2_b$ (fine-tuned MBERT) across three new
  MLLMs and the best existing model (GPT-4o mini$^\dagger$, from Table~\ref{tab:resultsSentVLMs}).
  \textbf{Bold}: best Task~1; \underline{underline}: best Task~$2_b$.%
}
\setlength{\tabcolsep}{5pt}
\renewcommand{\arraystretch}{1.25}
\centering

\begin{tabular}{c|cccc|cccc|cccc}
\hline
\multirow{2}{*}{\textbf{Problem}}
  & \multicolumn{4}{c|}{\textbf{Task 1 (Direct Classification)}}
  & \multicolumn{4}{c|}{\textbf{Task $2_a$ (MBERT pre-trained)}}
  & \multicolumn{4}{c}{\textbf{Task $2_b$ (MBERT fine-tuned)}}
\\
\cline{2-13}
  & \multicolumn{1}{c}{GPT-4o mini$^\dagger$}
  & \multicolumn{1}{c}{Gemini}
  & \multicolumn{1}{c}{Phi-4}
  & \multicolumn{1}{c|}{Gemma-4-E4B}
  & \multicolumn{1}{c}{GPT-4o mini$^\dagger$}
  & \multicolumn{1}{c}{Gemini}
  & \multicolumn{1}{c}{Phi-4}
  & \multicolumn{1}{c|}{Gemma-4-E4B}
  & \multicolumn{1}{c}{GPT-4o mini$^\dagger$}
  & \multicolumn{1}{c}{Gemini}
  & \multicolumn{1}{c}{Phi-4}
  & \multicolumn{1}{c}{Gemma-4-E4B}
\\
\hline
%% --- sigma3 P5 ---
$\langle\sigma_3,P_5\rangle$
  & $44.5$ \scriptsize $\pm 3.1$
  & $\mathbf{51.3}$ \scriptsize $\pm 1.5$
  & $42.2$ \scriptsize $\pm 3.7$
  & $46.8$ \scriptsize $\pm 4.8$
  & $47.9$ \scriptsize $\pm 2.3$
  & $48.3$ \scriptsize $\pm 2.3$
  & $22.5$ \scriptsize $\pm 24.3$
  & $31.3$ \scriptsize $\pm 29.1$
  & \underline{$58.4$} \scriptsize $\pm 4.0$
  & $56.9$ \scriptsize $\pm 1.3$
  & $53.2$ \scriptsize $\pm 2.8$
  & $57.4$ \scriptsize $\pm 2.8$
\\\hline
%% --- sigma3 P3 ---
$\langle\sigma_3,P_3\rangle$
  & $61.2$ \scriptsize $\pm 2.9$
  & $\mathbf{69.8}$ \scriptsize $\pm 1.3$
  & $61.2$ \scriptsize $\pm 2.9$
  & $60.7$ \scriptsize $\pm 2.8$
  & $66.6$ \scriptsize $\pm 6.1$
  & $69.9$ \scriptsize $\pm 3.1$
  & $50.6$ \scriptsize $\pm 31.4$
  & $67.3$ \scriptsize $\pm 4.0$
  & \underline{$77.5$} \scriptsize $\pm 0.9$
  & $76.3$ \scriptsize $\pm 3.0$
  & $73.1$ \scriptsize $\pm 2.7$
  & $75.4$ \scriptsize $\pm 1.2$
\\\hline
%% --- sigma5 P5 ---
$\langle\sigma_5,P_5\rangle$
  & $75.8$ \scriptsize $\pm 4.7$
  & $\mathbf{78.5}$ \scriptsize $\pm 6.0$
  & $78.1$ \scriptsize $\pm 6.7$
  & $78.3$ \scriptsize $\pm 4.2$
  & $72.2$ \scriptsize $\pm 4.4$
  & $70.4$ \scriptsize $\pm 10.7$
  & $14.4$ \scriptsize $\pm 32.6$
  & $11.0$ \scriptsize $\pm 20.7$
  & \underline{$84.4$} \scriptsize $\pm 4.2$
  & $81.4$ \scriptsize $\pm 6.6$
  & $79.0$ \scriptsize $\pm 5.5$
  & $83.3$ \scriptsize $\pm 4.8$
\\\hline
%% --- sigma5 P3 ---
$\langle\sigma_5,P_3\rangle$
  & $87.7$ \scriptsize $\pm 1.8$
  & $\mathbf{88.2}$ \scriptsize $\pm 2.2$
  & $82.1$ \scriptsize $\pm 1.5$
  & $81.0$ \scriptsize $\pm 1.6$
  & $90.6$ \scriptsize $\pm 1.5$
  & $88.3$ \scriptsize $\pm 1.4$
  & $87.8$ \scriptsize $\pm 3.2$
  & $77.9$ \scriptsize $\pm 26.7$
  & \underline{$95.8$} \scriptsize $\pm 0.9$
  & $94.5$ \scriptsize $\pm 2.1$
  & $93.3$ \scriptsize $\pm 2.1$
  & $94.2$ \scriptsize $\pm 0.7$
\\\hline
\end{tabular}
\label{tab:extendedMLLMbenchmark}
\end{adjustwidth}
\end{table*}


%% ======================================================================
%% REVIEWER RESPONSE LETTER
%% ======================================================================

\bigskip
\hrule
\bigskip

\noindent\textbf{Response to Reviewer Comment -- Limited Benchmark Comparisons}

\medskip

We thank the reviewer for this suggestion and have extended the evaluation to include
three additional state-of-the-art MLLMs spanning complementary paradigms:
\textbf{Gemini} (Google, proprietary API), \textbf{Phi-4-reasoning-vision-15B}
(Microsoft, open-weights decoder), and \textbf{Gemma-4-E4B-it}
(Google DeepMind, open-weights instruction-tuned).
Results are reported in Table~\ref{tab:extendedMLLMbenchmark}. On Task~1 (direct
classification), Gemini achieves the highest $F$-scores among all models tested
($51.3$--$88.2$\%), while Phi-4 and Gemma-4-E4B match GPT-4o mini under high-agreement
settings but lag behind on the harder low-agreement, five-class configuration
($42.2$\% and $46.8$\%, respectively, vs.\ $44.5$\% for GPT-4o mini).

Extending the evaluation to Task~$2_b$ (fine-tuned ModernBERT on MLLM-generated
descriptions) reveals that the description-mediated pipeline consistently outperforms
direct Task~1 classification for every newly added model:
absolute gains of $+2.9$--$+6.5$ pp for Gemini and $+0.9$--$+11.9$ pp for Phi-4.
Although GPT-4o mini + MBERT retains the highest absolute scores ($58.4$--$95.8$\%),
the margins over Gemini + MBERT and Phi-4 + MBERT are at most $3.5$ pp across all
settings, indicating that the performance advantage is structural --- attributable to
the pipeline architecture --- rather than to any particular MLLM choice.

Regarding the comparison with fusion-based multimodal models, we acknowledge that
our pipeline differs architecturally from early/late feature fusion approaches.
We clarify that the description-mediated design constitutes a form of
\emph{language-bridged modality fusion}: the MLLM externalises visual sentiment cues
as natural language, which the LLM then classifies, a strategy that is both more
interpretable and more composable than black-box feature concatenation.
We have added a paragraph to the Related Work (Section~2) that explicitly situates
our approach relative to representative fusion-based methods
\cite{PANDEY2024111206,das2023image,zhu2022multimodal} and discusses the
structural distinctions and the empirical advantages observed across our extended
set of models.
